\documentclass[11pt,a4paper]{article}

\usepackage[T1]{fontenc}
\usepackage[utf8]{inputenc}
\usepackage{amsmath,amssymb,amsthm}
\usepackage{mathtools}
\usepackage{hyperref}
\usepackage{doi}
\usepackage{geometry}
\usepackage{booktabs}
\usepackage{array}
\usepackage{enumitem}

\geometry{margin=2.5cm}

\hypersetup{
  colorlinks=true,
  linkcolor=black,
  citecolor=black,
  urlcolor=black
}

\theoremstyle{definition}
\newtheorem{remark}{Remark}

\newcommand{\SU}{\mathrm{SU}}
\newcommand{\Heis}{\mathrm{Heis}}
\newcommand{\ImH}{\operatorname{Im}\mathbb{H}}
\newcommand{\su}{\mathfrak{su}}
\newcommand{\Vr}{V_\rho}
\newcommand{\twoI}{2I}

\title{The Quantum Structure Sub-Programme\\[0.4em]
\large Presentation Note 7}

\author{J\'{e}r\^{o}me Beau\\[0.3em]
\small Independent Researcher\\
\small \texttt{jerome.beau@cosmochrony.org}}

\date{2026}

\begin{document}

  \maketitle

  \begin{abstract}
    The quantum structure sub-programme comprises three companion papers on the binary
    icosahedral group $\twoI$, plus one independently published result. \textbf{Q1} proves an
    exact, unconditional Fourier-support rigidity theorem in the finite Weil representation of
    the Heisenberg group, with a precise no-go corollary for support-, rank-, and
    Gram-based diagnostics; it establishes no phase coherence, singlet correlator, Tsirelson
    bound, or Born rule. \textbf{Q2} proves an exact finite-group fact: the spin-$\tfrac12$ and
    spin-$\tfrac32$ representations of $\twoI$ induce the same Cayley-graph Laplacian
    eigenvalue; it establishes no Casimir/isotropy normalisation, Born rule, Tsirelson-type
    bound, or $\SU(2)$ fixed-point/sector-selection argument. \textbf{Q3} proves a theorem
    conditional on an explicit, undemonstrated hypothesis: given that a bipartite state on a
    supplied composition is invariant under the diagonal action of $\twoI$ --- not derived from
    admissibility or from Born--Infeld indiscernibility --- the state is uniquely the $\SU(2)$
    singlet, with an explicit Casimir correlator, for all five admissible sectors of $\twoI$; it
    does not derive that invariance hypothesis, phase coherence, or the Born rule. No chain
    currently connects these three results to each other or to the structures of quantum
    mechanics. The \textbf{Bell paper}, published independently in \textit{Quantum Reports},
    establishes the structural non-applicability of Bell-type factorizability in non-injective
    effective descriptions; the proof of that central theorem uses only non-injectivity and
    informational completeness. This note records the current, narrow scope of each constituent and the
    open problems that remain.
  \end{abstract}

  \medskip
  \noindent\textbf{Keywords.}
  quantum structure; admissibility; binary icosahedral group; Fourier-support rigidity;
  conditional singlet theorem; Bell non-applicability; SU(2); non-injective projection;
  presentation note.

  \newpage
  \tableofcontents

  \newpage

  \section{Motivation and Current Status}
    \label{sec:motivation}

    This note surveys the quantum structure sub-programme's three constituent papers on the
    binary icosahedral group $\twoI$, plus the independently published Bell result, and states
    plainly what each does and does not establish.

    \medskip
    \noindent\textit{Original question.}
    The admissibility axioms A1--A4 force the fibre $F_n \simeq \Vr$ with Heisenberg
    structure (Note~5), and the admissibility thread $Q_8 \subset \twoI \subset \SU(2)$, with
    the spin-$\tfrac{1}{2}$ sector identified as the physically admissible sector (Note~1). The
    representation $\Vr$ already carries the complex $\SU(2)$ amplitude structure of the Weil
    representation, supplied by O18 and O23, not derived from A1--A4. \textit{Given that supplied
    carrier, do the structures of quantum mechanics built upon it --- phase coherence, the Born
    rule, a singlet correlator, Bell-type correlations --- follow from admissibility alone, or
    must they be separately postulated?}

    \medskip
    \noindent\textit{Current answer.}
    That derivation is not established. Q1's proved theorem is a representation-theoretic
    rigidity and no-go result, unrelated to phase coherence or the Born rule. Q2 proves only a
    finite Laplacian-eigenvalue coincidence. Q3 proves a theorem conditional on an explicit,
    undemonstrated invariance hypothesis, not derived from admissibility or from Born--Infeld
    indiscernibility. No result in this sub-programme derives phase coherence, a singlet
    correlator, a Tsirelson bound, the Born rule, or a physically preferred sector from the
    admissibility axioms alone.

    \begin{remark}[Vocabulary note]
      Throughout this note, admissibility replaces the earlier relaxation vocabulary. The
      substrate $\chi$ is static. Quantum structure, where established, is not the dynamics of
      $\chi$ but the algebraic consequence of non-injective admissible transitions.
    \end{remark}

    \begin{remark}[Scope]
      None of Q1, Q2, or Q3 derives phase coherence, the Born rule, a singlet correlator, or a
      Tsirelson-type bound from the admissibility axioms. Q1's rigidity theorem and Q2's
      eigenvalue degeneracy are unconditional but narrow. Q3's singlet theorem is conditional on
      an explicit, stated hypothesis that is itself not derived. The Bell paper's
      non-applicability result is unconditional; the proof of its central theorem uses only
      non-injectivity and informational completeness, independent of Q1, Q2, or Q3.
    \end{remark}

  \section{Position in the Cosmochrony Programme}
    \label{sec:position}

    The quantum structure sub-programme occupies an early and independent position in
    Branch~III. It depends on Branch~I (axioms, $F_n \simeq \Vr$, BI parity) and on selected
    O-series results (O18, O22, O23), but not on the geometry, gauge, or gravity sub-programmes.

    Of the results once listed here as upstream inputs --- the admissible fibre and Heisenberg
    structure (Note~5, Foundation~\cite{Beau2026m}, HeisenbergStructure~\cite{Beau2026HeisenbergStructure}),
    $\ImH \cong \su(2)$ and isotropic saturation (O23~\cite{Beau2026a27}), the parity involution
    and unit-norm amplitudes (O18~\cite{Beau2026a22}, O22~\cite{Beau2026a26}), and BI
    indiscernibility of conjugate Weil blocks (BornInfeld~\cite{Beau2026c}) --- none is actually
    used by the current, narrow content of Q1, Q2, or Q3. The only upstream input the current
    papers rely on is the Laplacian-eigenvalue data of the spectral admissibility programme
    (Note~1, SpectralAdmissibility~\cite{Beau2026a1}), which fixes which representation sectors of
    $\twoI$ are admissible and is used by Q2 and Q3.

    \noindent No output of this sub-programme currently supplies phase coherence, a correlator,
    a Tsirelson bound, an $\SU(2)$ fixed-point identification, or the Born rule to any downstream
    branch. Any paper citing this sub-programme for those results should be re-audited against
    the current, narrow status recorded below.

  \section{Constituent Papers}
    \label{sec:papers}

    \subsection{Q1 --- Fourier-support rigidity}
      \label{sec:q1}

      \textbf{Q1}~\cite{Beau2026q1} proves an exact, unconditional identity in the finite Weil
      representation of the Heisenberg group: for every complete BFS shell of the associated
      Cayley graph, conjugate Weil-sector fingerprint vectors span identical subspaces of
      $\mathbb{C}^q$. A corollary characterises precisely which diagnostics --- support, rank,
      subspace projectors, dimensions, principal angles, and any rephasing/permutation-invariant
      Gram-derived quantity --- are structurally blind to the conjugate-sector distinction.
      Verified independently by $1{,}548$ exact integer-arithmetic checks across six primes with
      zero failures, plus a negative control on a deliberately truncated shell.

      \textbf{What Q1 does not establish}: phase coherence between conjugate Weil blocks, a
      singlet correlator, a Tsirelson bound, or the Born rule. Q1's own interpretive content is a
      cautionary methodological point: a numerical near-zero-residual observation can be forced
      by a construction's own structural degeneracy rather than signal an intended physical
      mechanism.

    \subsection{Q2 --- An exact eigenvalue degeneracy on $\twoI$}
      \label{sec:q2}

      \textbf{Q2}~\cite{Beau2026q2} proves an exact finite-group fact: on the binary icosahedral
      group $\twoI$ with the canonical generating set $S = 10a \cup 10b$, the spin-$\tfrac12$ and
      spin-$\tfrac32$ representations induce the same Cayley-graph Laplacian eigenvalue,
      $\lambda_{1/2} = \lambda_{3/2} = 18$, proved via Schur's lemma on the central group-algebra
      element of the generating-set union. Under the admissibility-window definition
      $A^{\max}_\rho = c_{\mathrm{BI}}/\sqrt{\lambda_\rho}$, the two sectors have identical
      admissibility windows on $\twoI$.

      \textbf{What Q2 does not establish}: phase coherence, a Casimir or isotropy normalisation,
      a singlet correlator, a Tsirelson-type bound, the Born rule, or an $\SU(2)$
      fixed-point/sector-selection argument in any continuum or large-prime limit. None of these
      follow from an eigenvalue coincidence alone.

    \subsection{Q3 --- A conditional universal singlet}
      \label{sec:q3}

      \textbf{Q3}~\cite{Beau2026q3} proves a theorem conditional on an explicit hypothesis,
      stated as an assumption and not derived: for a supplied bipartite carrier
      $V_{\chi_{2j+1}} \otimes V_{\chi_{2j+1}}$, if a state on that carrier is invariant under the
      diagonal action of $\twoI$ (hypothesis \textbf{(H-inv)}), then the state is uniquely the
      $\SU(2)$ singlet $|\Omega_j\rangle$ up to phase, for each of the five admissible sectors
      $j \in \{\tfrac{1}{2}, 1, \tfrac{3}{2}, 2, \tfrac{5}{2}\}$ of $\twoI$. The proof uses
      $(V_j \otimes V_j)^{\twoI} \cong \mathrm{Hom}_{\twoI}(V_j^*, V_j)$, one-dimensional by
      Schur's lemma given irreducibility (an instance of the classical McKay correspondence for
      $\twoI$) and self-duality (every $\SU(2)$ character is real) of the restricted module.
      Given the singlet, the
      two-point correlator follows by an explicit trace computation:
      $E(\hat{a},\hat{b}) = -\tfrac{j(j+1)}{3}(\hat{a}\cdot\hat{b})$.

      \textbf{What Q3 does not establish}: (H-inv) itself --- it is not derived from
      admissibility or from any Born--Infeld indiscernibility argument, and is stated as a bare
      hypothesis throughout. Q3 also does not derive phase coherence, the Born rule, or a
      CHSH/Tsirelson-type bound, and does not select any sector as physically preferred.

    \subsection{Bell paper --- Bell non-applicability}
      \label{sec:bell}

      \textbf{Bell paper}~\cite{Beau2026b} is published in \textit{Quantum Reports}
      (MDPI, 2026) and is the only published result of the sub-programme. It establishes that
      Bell-type factorizability does not apply in non-injective effective descriptions, and does
      not depend on Q1, Q2, or Q3.

      \textit{Non-applicability of Bell factorizability} (proved). In any framework where
      observable outcomes arise as equivalence classes of underlying configurations under a
      non-injective map, the standard probabilistic factorization assumption implicit in Bell's
      theorem does not hold. The proof is structural: it requires only the non-injectivity of
      $\Pi$ and the informational completeness condition (E2 of ENI~\cite{Beau2026l}). No hidden
      variable, no non-local dynamics, and no modification of quantum mechanics is introduced.

      \textit{Mechanism}: observable outcomes correspond to equivalence classes
      $\Pi^{-1}(O_n)$; distinct configurations within the same class cannot be assigned
      independent local variables, preventing the joint factorization $p(a,b|\hat{a},\hat{b}) =
  p(a|\hat{a},\lambda) \cdot p(b|\hat{b},\lambda)$ for any common $\lambda$.

      \textit{Toy model}: an explicit model produces Bell--CHSH violations while preserving
      operational no-signalling and statistical independence of measurement settings. The model
      may exceed the Tsirelson bound (its role is to isolate the structural mechanism, not to
      reproduce quantum correlations quantitatively).

      \textit{Operator reformulation}: in the operator language, the same mechanism reformulates
      as a reduction of the observable algebra to a non-commutative conditional expectation.

      \begin{table}[ht]
        \centering
        \caption{Summary of results by paper. Status: P~=~proved (unconditional);
        C~=~conditional on an explicitly stated, undemonstrated hypothesis; N~=~numerical
        evidence.}
        \label{tab:papers}
        \renewcommand{\arraystretch}{1.2}
        \begin{tabular}{@{}lll@{}}
          \toprule
          Paper & Central output & Status \\
          \midrule
          Q1 & Fourier-support rigidity theorem and no-go corollary & P \\
          Q1 numerical & $1{,}548$ exact checks, six primes, zero failures & N \\
          Q2 & $\lambda_{1/2} = \lambda_{3/2} = 18$ on $\twoI$ (eigenvalue degeneracy) & P \\
          Q3 & Singlet $|\Omega_j\rangle$ and Casimir correlator, given (H-inv) & C \\
          Bell & Bell factorizability inapplicable in non-injective frameworks & P \\
          \bottomrule
        \end{tabular}
      \end{table}

  \section{Status of Results}
    \label{sec:status}

    \subsection{Proved (unconditional)}
      \label{sec:proved}

      \begin{enumerate}
        \item \textit{Fourier-support rigidity} (Q1): conjugate Weil-sector fingerprint vectors
        span identical subspaces on every complete BFS shell, with a precisely scoped no-go
        corollary.
        \item \textit{Eigenvalue degeneracy on $\twoI$} (Q2): $\lambda_{1/2} = \lambda_{3/2} = 18$
        for the canonical generating set.
        \item \textit{Bell factorizability non-applicable} (Bell paper~\cite{Beau2026b}): proved
        from non-injectivity and informational completeness of $\Pi$; published in
        \textit{Quantum Reports}.
      \end{enumerate}

    \subsection{Conditional}
      \label{sec:conditional}

      \begin{enumerate}
        \item \textit{Unique invariant singlet, all $j$} (Q3): given the explicit, undemonstrated
        hypothesis (H-inv) that a bipartite state on a supplied carrier is diagonally
        $\twoI$-invariant, the state is uniquely $|\Omega_j\rangle$ for each of the five
        admissible sectors.
        \item \textit{Casimir correlator} (Q3): given the singlet of the previous item,
        $E(\hat{a},\hat{b}) = -\tfrac{j(j+1)}{3}(\hat{a}\cdot\hat{b})$.
      \end{enumerate}

    \subsection{Numerical evidence}
      \label{sec:numerical}

      \begin{enumerate}
        \item \textit{Multi-prime exact verification} (Q1): $1{,}548$ (shell, block) checks
        across six primes, zero failures, plus a negative control confirming the criterion
        correctly detects failure on a deliberately truncated shell.
      \end{enumerate}

    \subsection{Open}
      \label{sec:open-hyp}

      \begin{enumerate}
        \item \textit{Phase coherence.} No result in this sub-programme derives phase coherence
        of the supplied Weil carrier from the admissibility axioms.
        \item \textit{The invariance hypothesis (H-inv).} Whether Q3's diagonal-invariance
        hypothesis follows from some further, separately stated admissibility or dynamical
        principle, rather than being supplied as an assumption, is open.
        \item \textit{The Born rule.} Not derived by Q1, Q2, or Q3, in the $\SU(2)$ sector or
        beyond it.
        \item \textit{A Tsirelson-type bound.} Not established by any constituent of this
        sub-programme.
        \item \textit{Sector selection.} No admissibility-flow or fixed-point argument
        currently selects a physically preferred sector among the admissible representations of
        $\twoI$.
        \item \textit{Multipartite entanglement.} The eigenvalue degeneracy of Q2 raises the
        question of entangled states mixing the spin-$\tfrac12$ and spin-$\tfrac32$ sectors; a
        multipartite extension beyond Q3's single-pair construction is open.
      \end{enumerate}

  \section{Conclusion}
    \label{sec:conclusion}

    The quantum structure sub-programme's three constituent papers on $\twoI$ establish, at
    present: an unconditional representation-theoretic rigidity theorem unrelated to quantum
    structure (Q1); an unconditional finite-group eigenvalue degeneracy (Q2); and a theorem
    conditional on an explicit, undemonstrated invariance hypothesis (Q3). No chain from
    admissibility, or from Born--Infeld indiscernibility, to phase coherence, a singlet
    correlator, a Tsirelson bound, the Born rule, or a physically preferred sector is currently
    established by these three papers. The independently published Bell paper establishes the
    structural non-applicability of Bell-type factorizability in non-injective effective
    descriptions; the proof of its central theorem uses only non-injectivity and informational
    completeness, independent of Q1, Q2, or Q3.

    The primary open problems are: deriving phase coherence of the supplied Weil carrier from
    the admissibility axioms; deriving Q3's invariance hypothesis (H-inv) rather than supplying
    it; and deriving the Born rule, in the $\SU(2)$ sector or beyond it.

    \newpage
    \bibliographystyle{plain}
    \bibliography{cosmochrony-bibliography}

\end{document}
